\documentclass{article}
\usepackage{../math_template}

\author{Jonathan Kasongo}
\title{Solutions to Advanced Problems by Stephen Siklos}

\begin{document}
\maketitle

\begin{problem}{Problem 1}
\begin{enumerate}
  \item[(i)] Find all sets of positive integers $a$, $b$ and $c$ that satisfy the equation
  \[
    \frac{1}{a} + \frac{1}{b} + \frac{1}{c} = 1.
  \]

  \item[(ii)] Determine the sets of positive integers $a$, $b$ and $c$ that satisfy the inequality
  \[
    \frac{1}{a} + \frac{1}{b} + \frac{1}{c} \geq 1.
  \]
\end{enumerate}
\end{problem}

\begin{solution}{Solution to Problem 1}
For part (i), the equation is symmetric so, without loss of generality,
assume that $c \leq b \leq a$. Multiply through by $abc$, to get
$ab + bc + ca = abc$. We claim that $c \leq 3$. Here is the proof,
Suppose, for the sake of contradiction that $c > 3$, then
$3ab < abc = ab + bc + ca \leq ab + ac + ab$ and so $ab < ac \iff b < c$
which is a contradiction since we assumed $b \geq c$. Now we can just check
each case.\\

\textit{Case 1:} $c = 1$. Then $ab + b + a = ab \iff a + b = 0$ but since
$a,b \geq 1$ this is impossible.\\

\textit{Case 2:} $c = 2$. Then $ab + 2(a+b) = 2ab \iff 2(a+b) = ab \iff
2b = ab - 2a \iff a = \frac{2b}{b-2}$. Since $a$ is a positive integer we
need $2b \equiv (b-2) + (b-2) + 4 \equiv 4 \equiv 0 \ (\text{mod} \ b-2)$.
Thus $b-2$ must be a factor of 4. We must have $b-2 = 1,2,4$ and so
$b = 3,4,6$. When $b = 6$, $a = \frac{12}{4} = 3$ which is impossible
since $a \geq b$. When $b = 4$, $a = \frac{8}{2} = 4$. When $b = 3$,
$a = \frac{6}{1}$. So the solutions here are $(a,b,c) = (4,4,2), (6,3,2)$.
\\

\textit{Case 3:} $c = 3$. Then $ab + 3(a + b) = 3ab \iff 3(a+b) = 2ab \iff
3b = 2ab - 3a \iff a = \frac{3b}{2b-3}$. Since $a$ is a positive integer
we need $3b \equiv b + 3 \equiv 0 \ (\text{mod} \ 2b-3)$. Notice that if
$2b - 3 > b + 3 \iff b > 6$, then we must have $b + 3 = 0 \iff b = -3$
since it can't be a positive multiple of $2b + 3$. But since $b$ is
positive this is impossible, so $3 = c \leq b \leq 6$. Now just try each
value of $b$. When $b = 3$, $a = \frac{9}{3} = 3$. When $b = 4$,
$a = \frac{12}{5} \notin \mathbb{N}$, so this is impossible. When $b = 5$,
$a = \frac{15}{7} \notin \mathbb{N}$, so this is impossible. When $b = 6$,
$a = \frac{18}{9} = 2$ but since $a \geq b$ this is impossible. So the only
solution here is $(a,b,c) = (3,3,3)$.\\

Finally combining each of the three cases shows that the only solutions are
$(a,b,c) = (4,4,2), (6,3,2)$ and $(3,3,3)$ up to permutation.\\

Part (ii) follows similarly, suppose that $c \geq 3$, then
$\frac{1}{a} \leq \frac{1}{b} \leq \frac{1}{c} \leq \frac{1}{3}$ but this
forces the equality case $a=b=c=3$. Now what if $c = 2$, then
$\frac{1}{a} + \frac{1}{b} \geq \frac{1}{2}$. We must have $b \leq 4$,
because if $b > 5$ then $\frac{1}{5} > \frac{1}{b} \geq \frac{1}{a}$ and
there is no way that they can sum to more than $\frac{1}{2}$. When $b = 2$,
$a$ can be any natural number. When $b = 3$, $\frac{1}{a} \geq \frac{1}{6}
\iff a\leq 6$. Now what if $c=1$, then $a,b$ can be any natural numbers.\\
\end{solution}

\end{document}
